% Transcription — fonds Alexandre Grothendieck, shelfmark 83, batch 2 (pages 21-40).
% Established from the facsimile archives/batches/83/batch-02.pdf (a mirror of
% https://grothendieck.umontpellier.fr/83.pdf), per the skill
% transcribe-grothendieck.
%
% Pass: Opus 5.5 (claude-opus-5-5), 2026-10-10 — first pass, unchecked against the pages by a human.
%
% Nineteen of the twenty pages are transcribed. Page 37 is a cover sheet
% inscribed in his hand ("Rapport anharmonique, avril 1986"); it has no
% \page{} and its words head the last section.
%
% The batch falls into three runs:
%   Pages 21-28 - scratch sheets (several written in both directions) of
%     matrix computations in SL_2 over a ring: e(alpha), e'(beta), u(lambda),
%     N(c), the Bruhat-type decompositions U_0 T U_1, N^- B_1, the relation
%     e'(b) e(-1/b) e'(b) = N(b); interleaved with fragments of another run
%     (author's pp. 114, 135; numbered relations (9)-(13), (64'), (70)) on a
%     group G acting on S, the double cover Z~ = G/G_0^+ of Z, and a
%     subgroup diagram, all running in from before the batch.
%   Pages 29-36 - a fair copy, author's pp. 1-4, "Structure du groupe
%     SL(2, Lambda)": split tori, U_0, U_1, B_0, B_1, N^-, sigma_0,
%     transporters, (sigma_0(b) e_0(b))^3 = 1, and the braid-type relation
%     u_0 u_1 u_0 = u_1 u_0 u_1 iff u_0 (x) u_1 = -1. Stops at (41).
%   Pages 38-40 - "A propos rapport anharmonique", April 1986, author's
%     pp. 1-3: three arrows V -> V_i in an abelian category, D, K(D, I),
%     the non-commutative variant with bitorsors, then four indices; it
%     runs out past the batch.
%
% Apparatus: 50 \ill{}, 39 \uncertain{}.
%
\documentclass[11pt,a4paper]{article}
\input{../preamble/grothendieck}

\folder{83}
\batch{2}
\pages{21}{40}

\dating{1982-1986}
\watermark{Édition de démonstration}
\foldertitle{[Icosaèdres, bi-icosaèdres et algèbres d'Azumaya de rang 4] : notes manuscrites (1982-1983, 1986, s.d.).}

\begin{document}

\section{Calculs dans $\mathrm{SL}_2$}
\note{Titre de l'éditeur. Les pages 21 et 22 sont des feuilles de calcul ; le calcul y court depuis les pages qui précèdent le lot (les éléments $e_0(\beta)$, $e_1(\gamma)$, $\sigma_0(\beta)$, $u(\lambda)$ sont supposés connus).}

\page{21}
\note{Feuille de calculs matriciels, sans texte suivi. Le symbole que nous transcrivons $\sigma_0$ est tracé comme un $v$ ou un $\nu$ à indice $0$ ; les pages 32 à 35 montrent qu'il s'agit de l'application $\sigma_0$ du texte mis au propre.}
\[
\begin{pmatrix} 0 & \beta \\ -\beta^{-1} & 0 \end{pmatrix}
\begin{pmatrix} 1 & \beta \\ 0 & 1 \end{pmatrix}
\begin{pmatrix} 0 & \beta \\ -\beta^{-1} & 0 \end{pmatrix}
= \begin{pmatrix} -1 & 0 \\ \beta^{-1} & -1 \end{pmatrix}
\qquad \sigma_0(\beta)^{-1} =
\]
le produit des deux premiers facteurs étant marqué d'une accolade :
\[
\begin{pmatrix} 0 & \beta \\ -\beta^{-1} & -1 \end{pmatrix}.
\]
\[
g = \begin{pmatrix} a & b \\ c & d \end{pmatrix}, \qquad
g^{-1} = \begin{pmatrix} d & -b \\ -c & a \end{pmatrix}
\]
\[
g \begin{pmatrix} 1 & \beta \\ 0 & 1 \end{pmatrix} g^{-1}
= \begin{pmatrix} a & a\beta + b \\ c & c\beta + d \end{pmatrix}
\begin{pmatrix} d & -b \\ -c & a \end{pmatrix}
= \begin{pmatrix} 1 - ac\beta & a^2\beta \\ -c^2\beta & 1 + ac\beta \end{pmatrix}
\]
\[
1 - a^2c^2\beta^2 \qquad\qquad ad - bc
\]
\struck{$e_1(\beta)$.} \quad \struck{$e_1(\gamma)\, e_0(-\gamma^{-1})\, e_1(\gamma) =$}
\[
\underbrace{\sigma_0(\beta)\, e_0(\beta)\, \sigma_0(\beta)^{-1}}\,
\underbrace{e_0(\beta)\, \sigma_0(\beta)\, e_0(\beta)\, \sigma_0(\beta)^{-1}} = \sigma_0(\beta)
\]
\note{Le dernier $\sigma_0(\beta)^{-1}$ porte sur le $\beta$ une surcharge que nous ne lisons pas ; les accolades sont de lui.}
\[
\omega\,\bigl(\sigma_0(\beta)\, e_0(\beta)\bigr)^3 =
\]
\[
\underbrace{\begin{pmatrix} 1 & 0 \\ \gamma & 1 \end{pmatrix}
\begin{pmatrix} 1 & \beta \\ 0 & 1 \end{pmatrix}}_{\begin{pmatrix} 1 & \beta \\ \gamma & \beta\gamma + 1 \end{pmatrix}}
\begin{pmatrix} 1 & 0 \\ \gamma & 1 \end{pmatrix}
= \begin{pmatrix} 1 + \beta\gamma & \beta \\ 2\gamma + \beta\gamma^2 & 1 + \beta\gamma \end{pmatrix},
\qquad 2\gamma + \beta\gamma^2 = \gamma(2 + \beta\gamma)
\]
\[
1 + \beta^2\gamma^2 + 2\beta\gamma - 2\beta\gamma - \beta^2\gamma^2
\]
\marginal{Écrit le long du bord gauche, la feuille tournée :
\[
\rho(\beta) = \begin{pmatrix} 0 & \beta \\ -\beta^{-1} & 0 \end{pmatrix}
\begin{pmatrix} 1 & \beta \\ 0 & 1 \end{pmatrix}
= \begin{pmatrix} 0 & \beta \\ -\beta^{-1} & -1 \end{pmatrix} = u(\beta)\,(
\]
\[
\begin{pmatrix} 0 & 1 \\ -1 & -1 \end{pmatrix}
\begin{pmatrix} 0 & 1 \\ -1 & -1 \end{pmatrix}
= \begin{pmatrix} -1 & -1 \\ 1 & 0 \end{pmatrix}
\]
\[
\rho'^3 = \begin{pmatrix} -1 & -1 \\ 1 & 0 \end{pmatrix}
\begin{pmatrix} 0 & 1 \\ -1 & -1 \end{pmatrix}
= \begin{pmatrix} 1 & 0 \\ 0 & 1 \end{pmatrix}
\]}
\note{La première ligne s'arrête sur « $u(\beta)\,($ », la parenthèse ouverte et vide. À la dernière, seule la première ligne du produit, « $1\ \ 0$ », est nettement lisible.}

\page{22}
\note{La page est écrite dans les deux sens. Une longue diagonale fine la traverse de part en part. En haut, dans le sens du papier :}
\[
\sigma_0(\beta)\bigl(e_0(\beta)\bigr) = e_1(-\beta^{-1})
\]
\[
\begin{pmatrix} 0 & \beta \\ -\beta^{-1} & 0 \end{pmatrix}
\begin{pmatrix} 1 & \beta \\ 0 & 1 \end{pmatrix} \bigl(
\]
\note{Le reste de la page est écrit tête-bêche ; nous le donnons dans l'ordre de lecture de la feuille retournée.}

\textbf{Remarque}\quad Considérons un multigraphe \struck{$(S, (A_\pi)_{\pi\in\Pi}, Z \subset$ \ill{}$)$}
\[
\mathcal{M} = \Bigl(\Gamma,\ Z \subset \mathrm{Aut}(\Gamma),\ \Pi \subset \Pi_Z \subset \mathrm{Bij}(S, \check S)\Bigr),
\qquad \Gamma = (S, \check S, R \subset S \times \check S)
\]
\note{Sous « $Z \subset \mathrm{Aut}(\Gamma)$ » : « sous-groupe » ; sous « $\Pi_Z$ » : « $Z$-torseurs ». Devant $\mathcal{M}$, un mot biffé illisible.}
avec op\supplied{ération} libre de $Z$ sur $S$ et $\check S$, il lui est donc associé

a) Un ens\supplied{emble} $T$ \ ($\simeq S/Z \simeq \check S/Z$)

b) \struck{Deux familles de} Un groupe $Z$ commutatif

c) Deux \struck{familles de} familles de $Z$-torseurs, \struck{\ill{}}, indexées par $T$ \struck{\ill{}} (savoir les fibres de $t$ dans $S$ et dans $\check S$ respectivement, i.e. les \uncertain{origines} et \uncertain{extrémités})

\ill{} \ill{}

\[
\begin{aligned}
&\chi\{s, t\} \in A, \qquad \chi\{s, \lambda t\} \in A, \qquad \chi\{\lambda s, t\} \in A, \qquad \chi\{s, \lambda^{-1} t\}
\end{aligned}
\]
\note{Colonne de droite, à côté d'un « $Z$ » isolé ; dans la marge, « $(s, \lambda t)$ ».}
\[
(s_0, \check s_0),\ (s_0, \lambda\check s_0) \in R; \qquad
\exists\, g \in G_0,\quad g(s_0) = s_0,\ g(\check s_0) = \lambda\check s_0
\]
\[
g^{-1}(s_0) = s_0, \qquad g^{-1}\lambda\, g^{-1}(s_0) =
\]
\[
g s_0 \qquad\qquad \lambda s_0 = s_0,\quad \lambda^{-1} t_0 = t_0
\]
$U g$ \qquad $t.u_0 =$ \struck{\ill{}} $t'u'_0 = \sigma(t'^{-1}u'_0)$, \qquad $T \in$
\struck{$t = v_0 u_0^{-1}$}

\struck{$u \in U_0$}\quad \struck{$u \in U_0 \cap$}\qquad $T \cap U_0$
\[
u s_0 = \lambda s_0 \qquad u t_0 = t_0
\]
\[
u \in T.U_0 \qquad u \in U_1
\]
\[
\boxed{TU_0 \cap U_1 = \{1\}}
\]
$TU_0 \cap U_1 = \{1\}$ (12) \quad \struck{$\Rightarrow$ Transitivité \ill{} $U_1$} \quad $\Longrightarrow \quad B_0 \cap B_1 = T$
\[
T \cap U_0 \subset U_1\,?
\]
\note{Les accolades de $\{1\}$ sont tracées comme un « $3$ » refermant le $1$. En marge gauche, la même égalité recopiée, et « $B_0 \cap B_1$ ».}

\page{23}
$M = L \oplus L'$ \qquad \struck{\ill{}} $L, L'$ $\Lambda$-modules inversibles
\[
\underbrace{\mathrm{SL}(M)}_{H} = \Bigl\{ \begin{pmatrix} a & b \\ c & d \end{pmatrix} \Bigm|
\begin{matrix} a, d \in \Lambda \\ c \in L^{-1}L',\ b \in L'^{-1}L \end{matrix},\ ad - bc = 1 \Bigr\}
\]
\note{Il écrit « $\mathrm{Se}(M)$ » ; nous lisons $\mathrm{SL}(M)$, et de même plus bas.}
\[
U_0 = \Bigl\{ \begin{pmatrix} 1 & \alpha \\ 0 & 1 \end{pmatrix} \Bigm| \alpha \in \overbrace{L'^{-1}L}^{L_0} \Bigr\}
\qquad
B_0 = \Bigl\{ \begin{pmatrix} \lambda & \alpha \\ 0 & \lambda^{-1} \end{pmatrix} \Bigm| \lambda \in \Lambda^*,\ \alpha \in L'^{-1}L \Bigr\}
\]
\[
U_1 = \Bigl\{ \begin{pmatrix} 1 & 0 \\ \beta & 1 \end{pmatrix} \Bigm| \beta \in \underbrace{L^{-1}L'}_{L_1} \Bigr\}
\qquad
B_1 = \Bigl\{ \begin{pmatrix} \lambda & 0 \\ \beta & \lambda^{-1} \end{pmatrix} \Bigm| \lambda \in \Lambda^*,\ \beta \in L^{-1}L' \Bigr\}
\]
\[
T = \Bigl\{ \begin{pmatrix} \lambda & 0 \\ 0 & \lambda^{-1} \end{pmatrix} \Bigm| \lambda \in \Lambda^* \Bigr\}
\qquad
\begin{pmatrix} a & b \\ c & d \end{pmatrix}^{-1} = \begin{pmatrix} d & -b \\ -c & a \end{pmatrix}
\]
\[
\sigma = \begin{pmatrix} 0 & -1 \\ 1 & 0 \end{pmatrix} \ \text{si } L = L' = \Lambda
\]
\struck{$\sigma(c)$ \ill{} $N(c) =$}
\[
N(c) = \begin{pmatrix} 0 & -c^{-1} \\ c & 0 \end{pmatrix} = \sigma . u(c)
\]
\[
g = \begin{pmatrix} a & b \\ c & d \end{pmatrix}, \quad
u(\lambda) = \begin{pmatrix} \lambda & 0 \\ 0 & \lambda^{-1} \end{pmatrix}, \quad
e(\alpha) = \begin{pmatrix} 1 & \alpha \\ 0 & 1 \end{pmatrix}, \quad
e'(\beta) = \begin{pmatrix} 1 & 0 \\ \beta & 1 \end{pmatrix}
\]
\[
N(c)\,u(\lambda) = N(c\lambda), \qquad u(\lambda)\,N(c) = N(\lambda^{-1}c), \qquad u(\lambda)\,N(c)\,u(\lambda') = N(\lambda^{-1}\lambda' c)
\]
\note{Le dernier argument est mal formé à l'encre ; la lecture $\lambda^{-1}\lambda' c$ est celle qu'imposent les deux égalités qui précèdent.}
Si $g = \begin{pmatrix} a & b \\ c & d \end{pmatrix}$, on a
\[
g \begin{pmatrix} \lambda & 0 \\ 0 & \lambda^{-1} \end{pmatrix} = \begin{pmatrix} a\lambda & b\lambda^{-1} \\ c\lambda & d\lambda^{-1} \end{pmatrix}, \qquad
\begin{pmatrix} \lambda & 0 \\ 0 & \lambda^{-1} \end{pmatrix} g = \begin{pmatrix} \lambda a & \lambda b \\ \lambda^{-1}c & \lambda^{-1}d \end{pmatrix}
\]
\[
\begin{aligned}
g \begin{pmatrix} \lambda & 0 \\ 0 & \lambda^{-1} \end{pmatrix} g^{-1}
&= \begin{pmatrix} a & b \\ c & d \end{pmatrix} \begin{pmatrix} \lambda d & -\lambda b \\ -\lambda^{-1}c & \lambda^{-1}a \end{pmatrix} \\
&= \begin{pmatrix} \lambda ad - \lambda^{-1}bc & ab(-\lambda + \lambda^{-1}) \\ cd(\lambda - \lambda^{-1}) & -\lambda bc + \lambda^{-1}ad \end{pmatrix}
\end{aligned}
\]
\[
N = \Bigl\{ \begin{pmatrix} a & b \\ c & d \end{pmatrix} \in \mathrm{SL}(M) \Bigm| ab = cd = 0 \Bigr\}
\]
\[
N^+ = T = \Bigl\{ \begin{pmatrix} a & b \\ c & d \end{pmatrix} \in \mathrm{SL}(M) \Bigm| b = c = 0 \Bigr\}, \qquad
N^- \,[= \sigma T] = \Bigl\{ \begin{pmatrix} a & b \\ c & d \end{pmatrix} \in \mathrm{SL}(M) \Bigm| a = d = 0 \Bigr\}
\]
\[
u(\lambda)\,\sigma = \begin{pmatrix} 0 & -\lambda \\ \lambda^{-1} & 0 \end{pmatrix}, \qquad
\sigma\, u(\lambda) = \begin{pmatrix} 0 & -\lambda^{-1} \\ \lambda & 0 \end{pmatrix}
\]
\note{Encadré dans le manuscrit :}
\[
e(\alpha)\,u(\lambda)\,e'(\beta) = \begin{pmatrix} 1 & \alpha \\ 0 & 1 \end{pmatrix}
\underbrace{\begin{pmatrix} \lambda & 0 \\ 0 & \lambda^{-1} \end{pmatrix} \begin{pmatrix} 1 & 0 \\ \beta & 1 \end{pmatrix}}_{\begin{pmatrix} \lambda & 0 \\ \lambda^{-1}\beta & \lambda^{-1} \end{pmatrix}}
= \begin{pmatrix} \lambda + \lambda^{-1}\alpha\beta & \lambda^{-1}\alpha \\ \lambda^{-1}\beta & \lambda^{-1} \end{pmatrix}
= \begin{pmatrix} a & b \\ c & d \end{pmatrix}
\]
\[
\lambda = d^{-1}, \qquad \alpha = \lambda b = b/d, \qquad \beta = \lambda c = c/d
\]
\note{Devant, un début biffé. Le « $\lambda b$ » est tracé comme $\lambda\beta$.}
donc
\[
U_0 \times T \times U_1 \longrightarrow H \quad \text{injectif}
\]
\[
U_0 T U_1 = \Bigl\{ \begin{pmatrix} a & b \\ c & d \end{pmatrix} \in H \Bigm| d \in \Lambda^* \Bigr\}
\]
\[
[e(\alpha), e'(\beta)] = \underbrace{e(\alpha)\,e'(\beta)}_{\begin{pmatrix} 1 + \alpha\beta & \alpha \\ \beta & 1 \end{pmatrix}}\, e(-\alpha)\,e'(-\beta)
\]
\[
\begin{pmatrix} 1 + \alpha\beta & \alpha \\ \beta & 1 \end{pmatrix}
\begin{pmatrix} 1 + \alpha\beta & -\alpha \\ -\beta & 1 \end{pmatrix}
= \begin{pmatrix} (1 + \alpha\beta)^2 - \alpha\beta & -\alpha^2\beta \\ \alpha\beta^2 & 1 - \alpha\beta \end{pmatrix}
= \begin{pmatrix} 1 + \alpha\beta + \alpha^2\beta^2 & -\alpha^2\beta \\ \alpha\beta^2 & 1 - \alpha\beta \end{pmatrix}
\]
si $\alpha\beta = 1$, on trouve
\[
[e(\alpha), e'(\beta)] = \begin{pmatrix} 3 & -\alpha \\ \alpha^{-1} & 0 \end{pmatrix}
\]
\struck{\ill{}}

donc si $\exists\,\lambda$ tel que $[e(\alpha), e'(\beta)] \in N^-$, alors $\alpha\beta = 1$, $3 = 0$ i.e. on est en car\supplied{actéristique} $3$.

\section{Le relèvement $\tilde\delta$ sur le centre}
\note{Titre de l'éditeur. La page 24 renvoie à une relation (64') et numérote (70) : la numérotation court depuis des pages antérieures au lot, dont $\tilde\delta$, $\tilde Z$, $u_\lambda$, $v_{\tilde\lambda}$, $T'$ et $G_0$ sont repris sans être redéfinis. Une longue diagonale fine traverse la page.}

\page{24}
La relation (64') donne, si $\lambda \in Z$, $\tilde\lambda \in \tilde Z$, avec $\lambda = \dot{\tilde\lambda} = \delta_{\tilde Z}(\tilde\lambda)$,
\struck{si $\lambda = \dot{\tilde\lambda}$, on a}
\[
(u_\lambda^{-1} v_{\tilde\lambda}^2)(s_0, t_0) = (\lambda s_0, \lambda t_0)
\]
\struck{\ill{} $u_\lambda^{-1} v_\lambda^2 (s_0, t_0$ \ill{}}

i.e.
\[
\lambda^{-1} u_\lambda^{-1} v_{\tilde\lambda}^2\, (s_0, t_0) = (s_0, t_0),
\]
ce qui implique que l'on a
\[
\lambda^{-1} u_\lambda^{-1} v_{\tilde\lambda}^2 \in T' \cap G_0 = \mathfrak{z}_{T'}
\]
\note{Le $\mathfrak{z}$ est la lettre en forme de « $3$ » bouclé dont il se sert pour un centre.}
mais on a aussi
\[
\tilde\delta(\lambda^{-1} u_\lambda^{-1} v_{\tilde\lambda}^2) = \tilde\delta(\lambda^{-1})\, \underbrace{\tilde\delta(u_\lambda^{-1})}_{1}\, \underbrace{\tilde\delta(v_{\tilde\lambda})^2}_{\tilde\lambda^2}
\]
je voudrais montrer que c'est $= 1$, i.e. qu'on a
\[
\tilde\delta(\lambda) = \tilde\lambda^2 \qquad \forall\, \tilde\lambda \in \tilde Z,\ \lambda = \dot{\tilde\lambda} \in Z \tag{70}
\]
i.e. qu'on ait \struck{\ill{}} commutativité dans

\begin{tikzcd}
\tilde Z \arrow[r, "{\delta_{\tilde Z} : \tilde\lambda \mapsto \dot{\tilde\lambda}}"] \arrow[rr, bend right=40, "{\tilde\lambda \mapsto \tilde\lambda^2}"'] & Z \arrow[r, "{\tilde\delta | Z}"] & \tilde Z
\end{tikzcd}

\section{Décompositions dans $H = \mathrm{SL}(M)$}
\note{Titre de l'éditeur. Les pages 23 et 25 à 28 reprennent les notations $e(\alpha)$, $e'(\beta)$, $u(\lambda)$, $N(c)$ de la page 23 ; les pages 26 à 28 sont encore des brouillons, l'exposé mis au propre commence page 29.}

\page{25}
\[
e'(\beta)\, e(\alpha) = \begin{pmatrix} 1 & 0 \\ \beta & 1 \end{pmatrix} \begin{pmatrix} 1 & \alpha \\ 0 & 1 \end{pmatrix}
= \begin{pmatrix} 1 & \alpha \\ \beta & 1 + \alpha\beta \end{pmatrix}
\]
si $1 + \alpha\beta \in \Lambda^*$, on a
\[
\boxed{\, e'(\beta)\, e(\alpha) = e\Bigl(\frac{\alpha}{1 + \alpha\beta}\Bigr)\, u\Bigl(\frac{1}{1 + \alpha\beta}\Bigr)\, e'\Bigl(\frac{\beta}{1 + \alpha\beta}\Bigr) \,}
\]
si $1 + \alpha\beta = 0$, i.e. $\alpha\beta = -1$, on a
\[
e'(\beta)\, e(\alpha) = \begin{pmatrix} 1 & \alpha \\ \beta & 0 \end{pmatrix} \overset{?}{\in} [\sigma B_1 =]\ N^- U_1
\]
\note{Suit un calcul biffé : $\sigma \begin{pmatrix} a' & \\ -\lambda^{-1} & \lambda^{-1} \end{pmatrix} \begin{pmatrix} 1 & 0 \\ \beta' & 1 \end{pmatrix} = \begin{pmatrix} -\lambda^{-1}\beta' & -\lambda^{-1} \\ \lambda & 0 \end{pmatrix} = \begin{pmatrix} a & b \\ c & d \end{pmatrix}$, avec en dessous $\begin{pmatrix} 0 & -1 \\ 1 & 0 \end{pmatrix} \begin{pmatrix} \lambda & 0 \\ \lambda^{-1}\beta' & \lambda^{-1} \end{pmatrix}$, barré lui aussi ; lecture incertaine dans le détail.}
\[
N^- B_1 = \Bigl\{ \begin{pmatrix} a & b \\ c & d \end{pmatrix} \in H \Bigm| d = 0 \Bigr\}
= \Bigl\{ \begin{pmatrix} a & b \\ c & 0 \end{pmatrix} \Bigm| a \in \Lambda,\ b \in L'^{-1}L,\ c \in L^{-1}L',\ bc = -1 \Bigr\}
\]
\[
\underbrace{\begin{pmatrix} 0 & b' \\ -b'^{-1} & 0 \end{pmatrix}}_{\in N^-}
\underbrace{\begin{pmatrix} 1 & 0 \\ \beta' & 1 \end{pmatrix}}_{\in U_1}
= \begin{pmatrix} b'\beta' & b' \\ -b'^{-1} & 0 \end{pmatrix}
\]
alors
\[
\begin{pmatrix} a & b \\ c & 0 \end{pmatrix} = \underbrace{\begin{pmatrix} 0 & b \\ c & 0 \end{pmatrix}}_{\in N^-} \underbrace{\begin{pmatrix} 1 & 0 \\ ab^{-1} & 1 \end{pmatrix}}_{\in U_1}
\qquad \text{si } bc = -1
\]
\struck{$e'(\beta)\, e(-\beta^{-1}) =$}

si $\alpha\beta = -1$
\[
e'(\beta)\, e(\alpha) = \underbrace{\begin{pmatrix} 0 & \alpha \\ \beta & 0 \end{pmatrix}}_{\in N^-} \begin{pmatrix} 1 & 0 \\ \alpha^{-1} = -\beta & 1 \end{pmatrix}
\qquad \text{i.e.}
\]
\[
e'(\beta)\, e(-\beta^{-1}) = \begin{pmatrix} 0 & -\beta^{-1} \\ \beta & 0 \end{pmatrix} \underbrace{\begin{pmatrix} 1 & 0 \\ -\beta & 1 \end{pmatrix}}_{e'(-\beta)}
\qquad \text{i.e.}
\]
\note{Sous la dernière matrice il écrit « $e'(\beta)$ », le signe moins absent ou indistinct.}
\[
e'(\beta)\, e(-\beta^{-1}) = N(\beta)\, e'(-\beta)
\]
\[
\boxed{\, e'(\beta)\, e(-\beta^{-1})\, e'(\beta) = N(\beta) \,} \qquad \beta \in (L^{-1}L')^* = U_1^*
\]
\struck{$e'(-\alpha^{-1})\, e(\alpha)\, e'(-\alpha^{-1}) = N(-\alpha^{-1})$}
\marginal{Dans la marge de droite, la feuille tournée, la vérification de l'encadré de la page 23 :
\[
\lambda + \lambda^{-1}\alpha\beta = d^{-1} + d\,\frac{bc}{d^2} = d^{-1} + \frac{bc}{d} = a
\]
i.e. \struck{\ill{}} OK}
\note{Au-dessus, dans la même marge, un calcul de puissances de $\alpha\beta$ ($1 - \alpha^3\beta^3 + \alpha^3\beta^3$…) rayé de traits ondulés ; au bas de la page, deux lignes écrites tête-bêche, barrées et rayées en croix, ne se lisent pas.}

\section{Le revêtement $\tilde Z$ et le diagramme des sous-groupes}
\note{Titre de l'éditeur. Toute la page 26 est écrite tête-bêche ; elle porte en tête le numéro « 114 », de sa main, et une diagonale fine la traverse. Les numéros (9) à (13) appartiennent à une suite dont le début est hors du lot.}

\page{26}
\note{Numéro de l'auteur : 114.}
On considère alors le groupe
\[
\tilde Z = G/G_0^+ \qquad \text{avec} \qquad G \xrightarrow{\ \tilde\delta\ } \tilde Z \tag{9}
\]
qui est un rev\supplied{êtement} de degré $2$ de $Z$
\[
1 \longrightarrow \{\pm 1\} \longrightarrow \tilde Z \longrightarrow Z \longrightarrow 1, \tag{10}
\]
et le diagramme d'inclusions de groupes

\begin{tikzcd}[column sep=large]
U \arrow[r, hook, "Z"] \arrow[d, hook, "\mu_2"'] & B'' \arrow[r, hook] \arrow[d, hook, "\mu_2"'] & G_0^+ \arrow[d, hook, "\mu_2"'] \\
\tilde U \arrow[r, hook, "Z"] \arrow[d, hook, "Z"'] & \tilde B'' \arrow[r, hook] \arrow[d, hook, "Z"'] & G_0 \arrow[d, hook, "Z"'] \\
B' \arrow[r, hook, "Z"] & B \arrow[r, hook] & G
\end{tikzcd}

\note{Diagramme (11). À gauche, une longue flèche verticale étiquetée $\tilde Z$ couvre les trois lignes ; à droite, de même, un trait ondulé marqué $\tilde Z$ entre deux flèches, l'une montante, l'autre descendante. Sous chacune des flèches horizontales marquées $Z$, une étiquette biffée.}
\marginal{On a écrit en dessous de toute flèche d'inclusion d'un sous-groupe distingué dans un groupe ambiant le groupe quotient correspondant (à isom\supplied{orphisme} près),}
défini alors à partir de
\[
G_0^+ \subset G_0 \subset G
\]
et de l'action de $G$ sur $S$, en termes d'un $G_0^+$-repère
\[
(s_0, t_0) \in \vec A_0 \qquad \text{i.e.} \quad s_0, t_0 \in S,\ \{s_0, t_0\} \in A_0 : \tag{12}
\]
\[
\left\{
\begin{aligned}
B' &= G_{s_0} = \{ g \in G \mid g s_0 = s_0 \} \\
\tilde U &= B' \cap G_0 = \{ g \in G_0 \mid g s_0 = s_0 \} = \{ b' \in B' \mid \delta(b') = 1 \} \\
U &= B' \cap G_0^+ = \{ g \in G_0^+ \mid g s_0 = s_0 \} = \{ b' \in B' \mid \tilde\delta(b') = 1 \} \\
\tilde B'' &= \{ g \in G_0 \mid g s_0 \in Z.s_0 \} \\
B'' &= \tilde B'' \cap G_0^+ \qquad [\mathrm{NB}\ \ \tilde B'' = \tilde U . B'']^{!} \\
B &= B' . B'' \ (= B' . \tilde B'')^{!}
\end{aligned}
\right. \tag{13}
\]
\marginal{En face de $\tilde B''$ : c'est le sous-groupe de $\mathrm{Norm}_{G_0}(\tilde U)$ ($\tilde U = G_{0\,s_0}$) qui \uncertain{sont} \ill{} \uncertain{décrivant} $Z \simeq \tilde B''/\tilde U$ \ill{} \uncertain{action} sur $S \simeq G_0/\tilde U$}
\note{Devant le crochet de la ligne de $B''$, un mot noirci ; à la ligne de $B$, une seconde copie « $\mathrm{NB}\ \tilde B'' = \tilde U . B''$ » est barrée, puis entourée et reliée par un trait au crochet de la ligne précédente.}
\marginal{Renvoi du signe « ! » : produits de sous-ens\supplied{embles} dans un groupe $G$ : $A.B = \{ ab \mid a \in A,\ b \in B \}$}
\marginal{En marge gauche : NB $B$ est bien un s\supplied{ous}-groupe car $B' = G_{s_0}$ normalise $B''$ qui est défini en termes de $s_0$}

\page{27}
\[
\bigl[e(\alpha)\, u(\lambda)\, \underbrace{e'(\beta)\bigr]\bigl[e(\alpha')}_{e\left(\frac{\alpha'}{1 + \alpha'\beta}\right) u\left(\frac{1}{1 + \alpha'\beta}\right) e'\left(\frac{\beta}{1 + \alpha'\beta}\right)}\, u(\lambda')\, e'(\beta')\bigr] =
\]
\[
\boxed{\, = e\Bigl(\alpha + \lambda^2 \frac{\alpha'}{1 + \alpha'\beta}\Bigr)\, u\Bigl(\frac{\lambda\lambda'}{1 + \alpha'\beta}\Bigr)\, e'\Bigl(\beta' + \lambda'^{-2} \frac{\beta}{1 + \alpha'\beta}\Bigr) \,}
\qquad 1 + \alpha'\beta \in \Lambda^*
\]
\note{Au-dessus, en ajout : « si $\alpha'\beta \neq -1$ », biffé, puis « $1 + \alpha'\beta \in \Lambda^*$ ».}
Cas $1 + \alpha'\beta = 0$, i.e. $\alpha'\beta = -1$ :
\[
e(\alpha)\, u(\lambda)\, \underbrace{e'(\beta)\, e(\alpha')}_{N(\beta)\, e'(-\beta)}\, u(\lambda')\, e'(\beta')
\]
\struck{\ill{}}
\[
\underbrace{N(\beta)\, u(\lambda')}_{\in N^-}\, \underbrace{e'(\beta' - \lambda'^{-2}\beta)}_{\in U_1}
= e(\alpha)\, u(\lambda)\, \underbrace{N(\beta)\, u(\lambda')}_{N(\beta\lambda^{-1}\lambda')}\, e'(\beta' - \lambda'^{-2}\beta)
\]
\[
e(\alpha)\, N(\underbrace{\beta\lambda^{-1}\lambda'}_{t})\, e'(\underbrace{\beta' - \lambda'^{-2}\beta}_{\gamma})
= N(t) . \underbrace{\underbrace{N(t)^{-1}(e(\alpha))}_{e'(-t^2\alpha)}\, e'(\gamma)}_{e'(-t^2\alpha + \gamma)}
\]
\[
e(\alpha)\, t\, e'(\gamma) = \begin{pmatrix} 1 & \alpha \\ 0 & 1 \end{pmatrix}
\underbrace{\begin{pmatrix} 0 & -t^{-1} \\ t & 0 \end{pmatrix} \begin{pmatrix} 1 & 0 \\ \gamma & 1 \end{pmatrix}}_{\begin{pmatrix} -t^{-1}\gamma & -t^{-1} \\ t & 0 \end{pmatrix}}
= \begin{pmatrix} -t^{-1}\gamma + \alpha t & -t^{-1} \\ t & 0 \end{pmatrix}
\]
\struck{$e(\alpha)\, N(t) = N(t)\, e'(\alpha')$}
\[
N(t)\bigl(e(\alpha)\bigr) \overset{\mathrm{def}}{=} N(t)\, e(\alpha)\, N(t)^{-1}
= \underbrace{\begin{pmatrix} 0 & -t^{-1} \\ t & 0 \end{pmatrix} \begin{pmatrix} 1 & \alpha \\ 0 & 1 \end{pmatrix}}_{\begin{pmatrix} 0 & -t^{-1} \\ t & t\alpha \end{pmatrix}}
\begin{pmatrix} 0 & t^{-1} \\ -t & 0 \end{pmatrix}
= \begin{pmatrix} 1 & 0 \\ -t^2\alpha & 1 \end{pmatrix}
\]
\[
\boxed{\, N(t)\, e(\alpha) = e'(-t^2\alpha) \,} \qquad t \in U_1^*,\ \alpha \in U_0
\]
\[
-t^2\alpha + \gamma = -\alpha\beta^2\lambda^{-2}\lambda'^2 + \beta' - \lambda'^{-2}\beta
\]
\[
e(\alpha)\, u(\lambda)\, e'(\beta)\, e(\alpha')\, u(\lambda')\, e'(\beta')
= \boxed{\, N(\beta\lambda^{-1}\lambda')\, e'\bigl(\beta' - \lambda'^{-2}\beta - \alpha\beta^2\lambda'^2\lambda^{-2}\bigr) \,}
\]
si $\alpha'\beta = -1$ i.e. $\alpha' = -\beta^{-1}$.

\page{28}
\note{La moitié supérieure de la page est un calcul dans le sens du papier ; la moitié inférieure, tête-bêche, numérotée « 135 » de sa main et traversée d'une diagonale, est un fragment de prose d'une autre suite.}
\[
\sigma \in B_G \qquad \text{\struck{$\sigma =$}}\ \sigma t \sigma t \qquad \sigma = t . u . t \qquad u\, g\, u^{-1} = g^{-1}
\]
\struck{$u g u^{-1} g^{-1} = g^{-2}$}
\[
\underbrace{\begin{pmatrix} 1 & 0 \\ \beta & 1 \end{pmatrix} \begin{pmatrix} 1 & -\beta^{-1} \\ 0 & 1 \end{pmatrix}}_{\begin{pmatrix} 1 & -\beta^{-1} \\ \beta & 0 \end{pmatrix}}
\begin{pmatrix} 1 & 0 \\ \beta & 1 \end{pmatrix}
\overset{?}{=} \begin{pmatrix} 0 & -\beta^{-1} \\ \beta & 0 \end{pmatrix} \qquad \text{OK}
\]
\note{« OK » est barré d'un trait, ou souligné en travers ; la lecture n'est pas sûre.}

\note{Partie tête-bêche, numéro de l'auteur 135 :}

Je m'aperçois : l'\uncertain{intérêt} que l'image des \add{sous-groupes} $T_0$ et de $T_1$ de $T_G$ dans $\mathrm{Aut}\, H$ sont un seul et \struck{\ill{}} sous-groupe (\add{et égal à l'image des $T_G$} — ce qui est évident aussi bien par les formules \struck{que}
\[
T_G = T_0 . T_H, \qquad \underbrace{\mathrm{Im}(T_H \to \mathrm{Aut}(H))}_{\subset\, \mathrm{Im}(T_0 \to \mathrm{Aut}(H))},
\]
qui impliquent que $T_0$ et $T_G$ ont \uncertain{donc} \struck{\ill{}} (et $T_1$ et $T_G$ aussi, par symétrie) ont \struck{\ill{}} \uncertain{images} dans $\mathrm{Aut}(H)$), le sous-groupe en question \uncertain{décidément} mérite d'être \uncertain{noté} $\tilde Z$, donc \struck{\ill{}} cette \uncertain{convention} d'écriture.

\[
\ill{} = \ill{}
\]
\note{Au bas de la partie retournée, deux formules courtes ($\ldots = \ldots u_\lambda^{-1}\ldots$), écrites dans le sens du papier, que nous ne lisons pas. Les deux lignes en tête de cette moitié, dans le sens du papier, se lisent : « $u_\lambda \sigma(s_0) = \lambda^{-1} t_0$ » et « $\sigma u_\lambda(s_0) = u_\lambda^{-1} \sigma(s_0) = \lambda t_0$ ».}

\section{Structure du groupe $\mathrm{SL}(2,\Lambda)$}
\note{Titre de l'auteur, souligné, en tête de la page 29, qui porte son numéro « 1 ». Les pages 29 à 40 forment une rédaction suivie, paginée par lui de 1 à 12 ; elle continue après le lot.}

\page{29}
\note{Numéro de l'auteur : 1.}
1. Soit $M$ un module libre de rang $2$ sur un anneau $\Lambda$, avec une structure symplectique
\[
e_M \in \bigl(\textstyle\bigwedge^2 M\bigr)^* . \tag{1}
\]
La donnée d'un tore maximal déployé $T$ de $\mathrm{SL}(M)$ revient à la donnée d'une décomposition de $M$ en somme de deux sous-modules libres \add{$L_i$} de rang $1$, formant un ensemble \struck{$\{L, L'\}$} à deux éléments $\varepsilon_T$ — \add{ce} \struck{\ill{}} sont les deux sous-modules propres de $M$, et on a
\[
\bigotimes_{i \in \varepsilon_T} L_i \simeq \bigl(\textstyle\bigwedge^2 M\bigr)_{\Lambda_{\pm 1}}^{\varepsilon_T} \simeq \Lambda_{\Lambda_{\pm 1}}^{\varepsilon_T} \tag{2}
\]
\note{Les indices et exposants de (2) sont serrés et la lecture du milieu, « $\Lambda_{\pm 1}$ », reste incertaine.}
en utilisant l'isomorphisme
\[
\textstyle\bigwedge^2 M \xrightarrow{\ \sim\ } \Lambda \qquad e_M \longmapsto 1 \tag{3}
\]
défini par la base $e_M$. L'opération de $T$ sur chaque $L_i$ définit un caractère
\[
T \xrightarrow{\ \chi_i\ } \mathbb{G}_m \qquad i \in \varepsilon_T \tag{4}
\]
qui est un isomorphisme, et (2) implique que l'on a
\[
\prod_{i \in \varepsilon_T} \chi_i = 1 \tag{5}
\]
i.e. les deux caractères $\chi_i$ sont inverses l'un de l'autre : ce sont les deux isomorphismes \add{déployés} de $T$ sur $\mathbb{G}_m$ (quand $\operatorname{Spec}\Lambda$ est connexe — sinon on peut dire \add{encore} que ce sont les deux iso\supplied{morphismes} « canoniques » de $T$ sur $\mathbb{G}_m$). Choisir l'un des \supplied{deux} isomorphismes revient donc à choisir
\marginal{Le long du bord gauche, en face des lignes 4 à 9, une note tournée, rattachée par un trait au mot « revient » souligné : « Si $\operatorname{Spec}\Lambda$ connexe, \ill{} \uncertain{sinon} \ill{} \uncertain{car} \ill{} \uncertain{lorsque} \ill{} \uncertain{donnée} \ill{} \uncertain{déployée} \ill{}, i.e. \ill{} $T \simeq \mathbb{G}_m$ \ill{} »}
\marginal{En bas à gauche, rattaché à « déployés » : « ils sont donc définis \uncertain{via} les \uncertain{dièdres} \ill{} »}

\page{30}
\note{Cette page ne porte pas de numéro de l'auteur. Il écrit désormais $L'_0$, $L'_1$ (avec un accent) pour les deux sous-modules propres.}
un des deux éléments de $\varepsilon_T$, i.e. à choisir un \emph{ordre} sur cet ensemble, qui nous permet dès lors de distinguer \add{un} $L_0$ et un $L_1$, de sorte que
\[
M = L'_0 \oplus L'_1 , \tag{6}
\]
et qu'on a un isomorphisme d'identification
\[
T \xrightarrow{\ \sim\ } \mathbb{G}_m, \tag{7}
\]
\add{un $\lambda \in \mathbb{G}_m(\cdot) \simeq T(\cdot)$} opérant sur $L_0$ par l'homothétie $\lambda$, sur $L_1$ par $\lambda^{-1}$.

On peut écrire, à l'aide de (6),
\[
H \overset{\mathrm{déf}}{=} \underbrace{\mathrm{SL}(M)}_{= \mathrm{Aut}(M, e_M)} \simeq
\Bigl\{ \begin{pmatrix} a & b \\ c & d \end{pmatrix} \Bigm|
\begin{matrix} a, d \in \Lambda,\ b \in \mathrm{Hom}(L'_1, L'_0) \simeq L_1'^{-1} L'_0 \\ c \in \mathrm{Hom}(L'_0, L'_1) \simeq L_0'^{-1} L'_1,\ ad - bc = 1 \end{matrix} \Bigr\} \tag{8}
\]
\note{Dans (8), « $\mathrm{SL}(M)$ » surcharge une première écriture, peut-être $\mathrm{SL}(2, M)(\Lambda)$.}
où
\[
bc = cb \in \Lambda \tag{9}
\]
désigne le composé de $b$ et $c$, identifié à un élément de $\Lambda$ (homothétie par $bc$).
\note{Dans (9), un premier « $\Lambda^*$ » est corrigé en $\Lambda$, l'astérisque restant au-dessus.}

La composition des éléments de $\mathrm{SL}(2, M)$ est donnée par la \emph{règle} usuelle. Avec ces notations, l'isomorphisme (7) est donné par
\[
\Lambda^* \xrightarrow{\ \sim\ } T(\Lambda) = \Bigl\{ \begin{pmatrix} \lambda & 0 \\ 0 & \lambda^{-1} \end{pmatrix} \Bigm| \lambda \in \Lambda^* \Bigr\} \tag{10}
\]
et on pose $u(\lambda) = \begin{pmatrix} \lambda & 0 \\ 0 & \lambda^{-1} \end{pmatrix}$ pour $\lambda \in \Lambda^*$.

Par la suite, on s'intéresse au groupe des points \struck{\ill{}} de $\mathrm{SL}(M)$ à valeurs dans $\Lambda$, plutôt qu'au schéma en groupes, et on notera par la suite par $T$ au lieu de $T(\Lambda)$, le sous-groupe

\page{31}
\note{Numéro de l'auteur : 2.}
(10) de $H$. On pose également
\[
L_0 = L_1'^{-1} L'_0, \qquad L_1 = L_0'^{-1} L'_1, \tag{11}
\]
ce sont donc des modules inversibles, canoniquement inverses l'un de l'autre. On a des homomorphismes canoniques
\[
e_0 : L_0 \longrightarrow H, \quad \beta \longmapsto \begin{pmatrix} 1 & \beta \\ 0 & 1 \end{pmatrix}, \qquad
e_1 : L_1 \longrightarrow H, \quad \gamma \longmapsto \begin{pmatrix} 1 & 0 \\ \gamma & 1 \end{pmatrix} \tag{12}
\]
et on pose
\[
\begin{aligned}
U_0 &= e_0(L_0) = \Bigl\{ \begin{pmatrix} 1 & \beta \\ 0 & 1 \end{pmatrix} \Bigm| \beta \in L_0 \Bigr\} \\
U_1 &= e_1(L_1) = \Bigl\{ \begin{pmatrix} 1 & 0 \\ \gamma & 1 \end{pmatrix} \Bigm| \gamma \in L_1 \Bigr\} .
\end{aligned} \tag{13}
\]
On a pour $\lambda \in \Lambda^*$, $g = \begin{pmatrix} a & b \\ c & d \end{pmatrix} \in H$,
\[
\left\{
\begin{aligned}
u(\lambda)\, g &= \begin{pmatrix} \lambda a & \lambda b \\ \lambda^{-1}c & \lambda^{-1}d \end{pmatrix} \qquad
g\, u(\lambda) = \begin{pmatrix} \lambda a & \lambda^{-1}b \\ \lambda c & \lambda^{-1}d \end{pmatrix} \\
u(\lambda)\, g\, u(\lambda)^{-1} &= \begin{pmatrix} a & \lambda^2 b \\ \lambda^{-2}c & d \end{pmatrix}
\end{aligned}
\right. \tag{14}
\]
en particulier
\[
\left\{
\begin{aligned}
u(\lambda)\bigl(e_0(\beta)\bigr) &= e_0(\lambda^2\beta) \\
u(\lambda)\bigl(e_1(\gamma)\bigr) &= e_0(\lambda^{-2}\gamma)
\end{aligned}
\right. \tag{15}
\]
\note{Sic : $e_0$ à la seconde ligne de (15), où l'on attend $e_1$.}
\marginal{attention à l'exposant $2$ des $\lambda$ dans la $2^{\text{e}}$ \uncertain{manière} !}
Ainsi $T$ normalise $U_0$ et $U_1$, et on a bien sûr
\[
T \cap U_0 = T \cap U_1 = \{1\}, \tag{16}
\]
on in\supplied{tro}duit
\[
\left\{
\begin{aligned}
B_0 = T . U_0 &= \{ e_0(\beta)\, u(\lambda) \mid \beta \in \Lambda,\ \lambda \in \Lambda^* \}
= \Bigl\{ \begin{pmatrix} \lambda & \lambda^{-1}\beta \\ 0 & \lambda^{-1} \end{pmatrix} \Bigm| \lambda \in \Lambda^*,\ \beta \in \Lambda \Bigr\} \\
&\qquad \Bigl( = \Bigl\{ \begin{pmatrix} a & b \\ c & d \end{pmatrix} \in H \Bigm| c = 0 \Bigr\} \Bigr) \\
B_1 = T . U_1 &= \{ u(\lambda)\, e_1(\gamma) \mid \gamma \in \Lambda,\ \lambda \in \Lambda^* \}
= \Bigl\{ \begin{pmatrix} \lambda & 0 \\ \lambda^{-1}\gamma & \lambda^{-1} \end{pmatrix} \Bigm| \lambda \in \Lambda^*,\ \gamma \in \Lambda \Bigr\} \\
&= \Bigl\{ \begin{pmatrix} a & b \\ c & d \end{pmatrix} \in H \Bigm| b = 0 \Bigr\}
\end{aligned}
\right. \tag{17}
\]
\note{Entre les deux lignes de (17), un « $B_1$ » commencé et abandonné.}
qui sont des sous-groupes de $H$ normalisant $U_0$ resp\supplied{ectivement} $U_1$. On a
\[
B_0 \cap B_1 = T = \Bigl\{ \begin{pmatrix} a & b \\ c & d \end{pmatrix} \in H \Bigm| b = c = 0 \Bigr\} . \tag{18}
\]

\page{32}
On introduit aussi une application (\emph{pas} un homomorphisme !)
\[
\left\{
\begin{aligned}
&\sigma_0 : L_0^* \longrightarrow H \\
&\sigma_0(\beta) = \begin{pmatrix} 0 & \beta \\ -\beta^{-1} & 0 \end{pmatrix}
\end{aligned}
\right. \tag{19}
\]
\note{Dans (19), le coefficient en haut à droite est surchargé ; nous lisons $\beta$, d'après (20) et les calculs de la page 21. Un signe d'égalité y est d'abord écrit puis biffé.}
\marginal{\emph{NB} Si $L_0$ n'est pas libre, alors $L_0^*$ est vide, et ce qui suit sans intérêt. On suppose donc par la suite $L_0$ libre, donc $L_1$ libre}
on a alors
\[
\left\{
\begin{aligned}
\sigma_0(\beta)^2 &= \begin{pmatrix} -1 & 0 \\ 0 & -1 \end{pmatrix} = u(-1) \overset{\mathrm{déf}}{=} \omega \qquad (\text{élément \emph{central} de } H) \\
\sigma_0(\beta)\bigl(u(\lambda)\bigr) &= u(\lambda^{-1}) \qquad \text{donc } \sigma(\beta) \text{ normalise } T \\
\sigma_0(\lambda\beta) &= u(\lambda)\, \sigma_0(\beta) = \sigma_0(\beta)\, u(\lambda^{-1}) ,
\end{aligned}
\right. \tag{20}
\]
\note{À la dernière ligne, un astérisque biffé après $\beta$ ; sous « $\sigma_0(\beta)\, u(\lambda^{-1})$ », deux mots biffés (« donc pour » ?).}
la dernière formule montre que, si on pose
\[
N^- = \sigma_0(L_0^*) = \Bigl\{ \begin{pmatrix} a & b \\ c & d \end{pmatrix} \in H \Bigm| a = 0,\ d = 0 \Bigr\}, \tag{21}
\]
\note{Devant $N^-$, un mot biffé ; avant l'accolade, « $H$ » biffé.}
alors $N^-$ est une classe à gauche ou à droite mod $T$, et que
\[
L_0^* \xrightarrow{\ \sim\ } N^- \tag{22}
\]
est bijection, et \struck{\ill{}} un isomorphisme de $\Lambda^*$-torseurs, quand on fait opérer $\Lambda^*$ sur $L_0^*$ par homothéties, et sur $N^-$ par multiplication \emph{à gauche} par $u(\lambda)$. On pourrait bien sûr introduire \struck{symétriquement le composé}
\[
\sigma_1 : L_1 \xrightarrow{\ \sim\ } N^- \qquad \sigma_1(\gamma) = \begin{pmatrix} 0 & -\gamma^{-1} \\ \gamma & 0 \end{pmatrix} \tag{23}
\]
de sorte que l'on a
\[
\left\{
\begin{aligned}
\sigma_1(\gamma) &= \sigma_0(-\gamma^{-1}) \\
\sigma_1(\lambda\gamma) &= u(\lambda^{-1})\, \sigma_1(\gamma) = \sigma_1(\gamma)\, u(\lambda)
\end{aligned}
\right. \tag{24}
\]
donc $\sigma_1$ est un isomorphisme de \add{$\Lambda^*$-}torseurs, quand cette fois on fait opérer $\Lambda^*$ sur $N^-$ par multiplication : \emph{droite} par $u(\lambda)$.

\page{33}
\note{Numéro de l'auteur : 3.}
Notons la formule fondamentale
\[
\boxed{\, \sigma_0(\beta)\bigl(e_0(\beta)\bigr) = e_1(-\beta^{-1}) \,} \tag{25}
\]
\note{Dans l'encadré, une matrice entièrement noircie précède $e_1(-\beta^{-1})$.}
\marginal{d'où plus généralement $\sigma_0(\beta)\bigl(e_0(\lambda\beta)\bigr) = e_1(-\lambda\beta^{-1})$}
qui montre en particulier que
\[
\forall\, \sigma \in N^-, \quad \text{on a} \quad \sigma(U_0) = U_1, \quad \text{d'où} \quad \sigma(U_1) = U_0 \tag{26}
\]
(puisque $\sigma^2 = \omega$ est central).

Notons que
\[
g = \begin{pmatrix} a & b \\ c & d \end{pmatrix} \in H \Longrightarrow g^{-1} = \begin{pmatrix} d & -b \\ -c & a \end{pmatrix} \tag{27}
\]
et
\[
g\, e_0(\beta)\, g^{-1} = \begin{pmatrix} 1 - ac\beta & a^2\beta \\ -c^2\beta & 1 + ac\beta \end{pmatrix} \tag{28}
\]
donc on a $g(U_0) \subset U_1$ ssi on a
\[
a^2\beta = 0, \quad ac\beta = 0 \qquad \forall\, \beta \in L_0
\]
i.e.
\[
a^2 = 0, \quad ac = 0
\]
d'où $a = a\underbrace{(ad - bc)}_{1} = 0$, donc
\[
\left\{
\begin{aligned}
\mathrm{Transp}_H(U_0, U_1) &= \Bigl\{ \begin{pmatrix} a & b \\ c & d \end{pmatrix} \in H \Bigm| a = 0 \Bigr\}, \quad \text{de même} \\
\mathrm{Transp}_H(U_1, U_0) &= \Bigl\{ \begin{pmatrix} a & b \\ c & d \end{pmatrix} \in H \Bigm| d = 0 \Bigr\} .
\end{aligned}
\right. \tag{29}
\]
\marginal{En marge de (29) : il suffit même qu'un $g \in H$ \struck{soit} tel que pour \emph{un} $\beta \in L_0^*$, on ait $g(\beta) \in L_1$, pour avoir $g \in \mathrm{Transp}_H(U_0, U_1)$}
\struck{je dis que}

En fait, \struck{les calculs montrent} que dans ces cas \uncertain{là}, le transporteur large est strict, car si $a = 0$, la formule (28) donne
\[
g\, e_0(\beta)\, g^{-1} = \begin{pmatrix} 1 & 0 \\ -c^2\beta & 1 \end{pmatrix} \tag{28'}
\]
or $ad - bc = -bc = 1$, donc $b$ et $c$ sont inversibles, donc $c^2$ inversible, donc $\mathrm{int}(g)$ induit un iso\supplied{morphisme} $U_0 \xrightarrow{\ \sim\ } U_1$. Des formules (29) on tire

\page{34}
\[
N^- = \mathrm{Transp}(U_1, U_0) \cap \mathrm{Transp}(U_0, U_1) \tag{30}
\]
Soient maintenant $\beta \in L_0$, $\gamma \in L_1$ avec
\[
\left\{
\begin{aligned}
&\beta\gamma = -1 \qquad \text{i.e. } \gamma = -\beta^{-1}, \text{ i.e. } \beta = -\gamma^{-1} \\
&\text{i.e. } e_1(\gamma) = \underbrace{\sigma_0(\beta)\bigl(e_0(\beta)\bigr)}_{e_1(-\beta^{-1})} ,
\end{aligned}
\right. \tag{31}
\]
je dis qu'on a alors
\[
e_1(\gamma)\, e_0(\beta)\, e_1(\gamma) = \sigma_0(\beta) \qquad (\text{pour } \beta\gamma = -1) \tag{32}
\]
formule qui s'écrit aussi, puisque $e_1(\gamma) = \sigma_0(\beta)\, e_0(\beta)\, \sigma_0(\beta)^{-1}$,
\[
\underbrace{\sigma_0(\beta)\, e_0(\beta)}\, \sigma_0(\beta)^{-1}\, e_0(\beta)\, \underbrace{\sigma_0(\beta)\, e_0(\beta)}\, \sigma_0(\beta)^{-1} = \sigma_0(\beta)
\]
\note{Le numéro de cette formule est surchargé, illisible.}
soit encore, en écrivant $\sigma_0(\beta)^{-1} = \omega\, \sigma_0(\beta)$ ($\omega$ central, $\omega^2 = 1$)
\[
\bigl(\sigma_0(\beta)\, e_0(\beta)\bigr)^3 = 1
\]
\note{Avant « $= 1$ », deux signes noircis.}
soit encore
\[
\rho(\beta)^3 = 1, \qquad \text{où } \rho(\beta) = \sigma_0(\beta)\, e_0(\beta), \tag{33}
\]
soit
\[
\rho(\beta) = \begin{pmatrix} 0 & \beta \\ -\beta^{-1} & -1 \end{pmatrix} = u(\beta)\, \rho' \qquad \rho' = \rho(1) = \begin{pmatrix} 0 & 1 \\ -1 & -1 \end{pmatrix} \tag{34}
\]
\note{Au-dessus de l'égalité, « soit $\rho(\beta)$ » biffé ; sous elle : « si $L_0 = L_1 = \Lambda$, donc $\beta \in \Lambda^*$ ».}
la vérification est immédiate.

Supposons maintenant $\beta \in L_0$, $\gamma \in L_1$ quelconques, et calculons $e_1(\gamma)\, e_0(\beta)\, e_1(\gamma)$, on trouve
\[
e_1(\gamma)\, e_0(\beta)\, e_1(\gamma) = \begin{pmatrix} 1 + \beta\gamma & \beta \\ \gamma(2 + \beta\gamma) & 1 + \beta\gamma \end{pmatrix} \tag{35}
\]
\note{Le numéro (35), comme (36) ensuite, est mal formé ; nous les supposons d'après la suite (34), (37).}
qui prouve
\[
\text{si } \beta \in L_0,\ \gamma \in L_1, \text{ les conditions suivantes sont équivalentes} \tag{37}
\]
\[
\left|
\begin{aligned}
&\text{a) } \beta\gamma = -1 \\
&\text{b) } e_1(\gamma)\, e_0(\beta)\, e_1(\gamma) \in \mathrm{Transp}(U_0, U_1) \\
&\text{c) } e_1(\gamma)\, e_0(\beta)\, e_1(\gamma) \in \mathrm{Transp}(U_1, U_0) \\
&\text{d) } e_1(\gamma)\, e_0(\beta)\, e_1(\gamma) \in N^- \\
&\text{e) } e_1(\gamma)\, e_0(\beta)\, e_1(\gamma) = \sigma_0(\beta)
\end{aligned}
\right.
\]
\marginal{En marge de (35) : faire le lien avec la formule pour $e_1(\beta)\, e_0(\alpha') = e_0(x)\, u(t)\, e_1(y)$ pour $1 + \beta\alpha'$ \ill{}.}

\page{35}
\note{Numéro de l'auteur : 4.}
En d'autres termes, en termes du groupe $H$, et des deux sous-groupes $U_0$, $U_1$ de $H$ et de leurs structures de $\Lambda$-modules inversibles $L_0$, $L_1$ (automatique si $\Lambda$ est un anneau \uncertain{égal} à $\mathbb{Z}$ ou $\mathbb{Z}/n\mathbb{Z}$), on trouve

a) L'accouplement $L_0 \otimes L_1 \simeq \Lambda$, par la condition que $u_0 \otimes u_1 = -1$ ssi $u_1 u_0 u_1 \in \mathrm{Transp}(U_0, U_1)$, ou encore $u_1 u_0 u_1 \in \mathrm{Transp}(U_0, U_1) \cap \mathrm{Transp}(U_1, U_0) \overset{\mathrm{déf}}{=} N^-$

b) L'application $\sigma_0 : U_0^* \longrightarrow N^-$ par la condition $\sigma_0(u_0) = u_1 u_0 u_1$, où $u_1 \in U_1$ est l'unique élément de $U_1$ satisfaisant à la condition que $u_1 u_0 u_1 \in N^-$, et alors on a

c) $u_1 = \sigma_0(u_0)(u_0) = \underbrace{u_1 u_0 u_1}\, u_0\, \underbrace{u_1^{-1} u_0^{-1} u_1^{-1}}$

\marginal{En face de a) et b), reliée par une accolade : ou encore $u_0 \in U_0^*$ et \struck{$(u_1 u_0 u_1)$} $\mathrm{int}(u_1 u_0 u_1)(u_0) = u_1$, i.e. (38) ci-dessous}
i.e.
\[
u_0 u_1 u_0 u_1^{-1} u_0^{-1} u_1^{-1} = 1 \qquad \text{i.e. } u_0 u_1 u_0 = u_1 u_0 u_1 \tag{38}
\]
formule qui est équivalente encore à $u_0 \otimes u_1 = -1$. En termes de $e_0$, $e_1$, $\beta\gamma$, s'écrit aussi
\[
e_0(\beta)\, e_1(\gamma)\, e_0(\beta) = e_1(\gamma)\, e_0(\beta)\, e_1(\gamma) \quad \text{ssi} \quad \beta\gamma = -1 \qquad (\beta \in L_0^*,\ \gamma \in L_1) . \tag{39}
\]
La relation (38) \struck{est} entre $U_0$ et $U_1$ est le graphe d'une bijection entre $U_0^* = e_0(L_0^*)$ et $U_1^* = e_1(L_1^*)$, qu'on peut noter
\[
u_0 \longmapsto u_0^\vee : U_0^* \longrightarrow U_1^*
\]
et satisfaisant
\[
(u_0^\lambda)^\vee = (u_0^\vee)^{\lambda^{-1}}, \qquad u_0 \in U_0^*,\ \lambda \in \Lambda^* \tag{40}
\]
\note{Dans (40), $u_0^\lambda$ surcharge une première écriture $\lambda u_0$.}
en notant exponentiellement l'opération \add{externe} de $\Lambda$ sur $U_0$, $U_1$.

\page{36}
(alors que pour $g \in \mathrm{Transp}_H(U_0, U_1)$, la bijection
\[
\mathrm{int}(g) | U_0 : U_0 \xrightarrow{\ \sim\ } U_1
\]
respecte les structures de $\Lambda$-module sur $U_0$ et $U_1$, cf. (28')).

\struck{On notera}
\[
N' = T \sqcup N^- \subset H \tag{41}
\]
\struck{est un sous-groupe de}
\note{La formule (41) est enfermée dans un cadre qui barre les mots au-dessus et au-dessous ; le reste de la page est blanc.}

\section{Rapport anharmonique}
\note{Titre de l'auteur, pris sur le feuillet qui sert de couverture à la suite (page 37 du fonds, non transcrite) : « Rapport anharmonique — avril 1986 ». La rédaction qui suit, datée « Avril 86 » en tête de la page 38, est paginée par lui à partir de 1 ; elle continue après le lot.}

\page{38}
\note{Numéro de l'auteur : 1. En haut à droite : « Avril 86 ».}
À propos rapport anharmonique.

(a) Soient $V \xrightarrow{\ \varphi_i\ } V_i$ ($i \in I$, $\operatorname{card} I = 3$) des flèches de même source dans une cat\supplied{égorie} abélienne $\mathcal{C}$, telles que $\forall J \subset I$, $\operatorname{card} J = 2$, la flèche correspondante \struck{soit}
\[
V \longrightarrow V_J \overset{\mathrm{déf}}{=} \prod_{i \in J} V_i
\]
soit iso\supplied{morphisme}. \struck{Considérons} Alors on a suite exacte
\[
0 \longrightarrow V \xrightarrow{\ \varphi\ } \prod V_i \longrightarrow D \longrightarrow 0, \qquad \varphi = (\varphi_i), \quad D = \prod V_i / \mathrm{Im}\, V \tag{1}
\]
et les homomorphismes composés
\[
u_i : V_i \xrightarrow{\ \mathrm{inj}_i\ } \prod V_i \longrightarrow D
\]
\note{Il dessine $u_i$ comme une flèche courbe sous la composée.}
sont iso\supplied{morphismes}. On a donc un iso\supplied{morphisme} des suites exactes entre (1) et
\[
0 \longrightarrow K(D, I) \longrightarrow D^I \xrightarrow{\ \varepsilon_{D,I}\ } D \longrightarrow 0
\]
où
\[
K(D, I) \overset{\mathrm{déf}}{=} \operatorname{Ker} \varepsilon_{D,I} = \Bigl\{ (x_i)_{i \in I} \in D^I \Bigm| \textstyle\sum x_i = 0 \Bigr\}
\]
\note{Sous $\varepsilon_{D,I}$, un mot souligné, peut-être « somme ».}
Ainsi, la catégorie des systèmes $(V, (V_i), (\varphi_i))$ \uncertain{est} équivalente à celle des paires $(D, I)$, i.e. est équivalente à : $\mathrm{Ens}_3 \times \mathcal{C}$.

(b) \emph{NB} Il y a une variante non commutative où $V$, $V_i$ sont des \uncertain{ensembles} \struck{avec \uncertain{dimensions} de} \ill{} plus. \struck{\ill{}} On doit pouvoir montrer

\page{39}
\note{Numéro de l'auteur : 2. Les indices sont tantôt $h$, tantôt $k$ pour le troisième élément de $I$ ; nous les gardons tels qu'ils sont écrits.}
(utilisant une orientation de $I$) qu'il existe sur chaque $V_i$ une structure de bitorseur \add{sous} deux paires de groupes
\[
G_j,\ G_h \qquad G_j = G_{\bar\rho(i)},\quad G_h = G_{\rho(i)}
\]
(\uncertain{les} $G_j$ étant définis \struck{\ill{}} à isom\supplied{orphisme} unique près) et un iso\supplied{morphisme} (choisissant un $i \in I$) par la commodité de l'écriture
\[
V_i \mathbin{\wedge_{G_k}} V_j \mathbin{\wedge_{G_i}} V_k \simeq \mathbb{1}_{G_i} \qquad \text{isom\supplied{orphisme} de bitorseurs}
\]
\note{Sous le premier $\wedge$, « $G_k$ » suivi d'un mot noirci.}
\drawing{En marge gauche, un triangle : sommets $V_h$ (en haut), $V_i$ et $V_j$ (en bas) ; côtés marqués $G_j$ (à gauche), $G_i$ (à droite), $G_k$ (en bas).}
i.e. un élément \add{$\varepsilon$} central \uncertain{dans} $V_i \mathbin{\wedge_{G_k}} V_j \mathbin{\wedge_{G_i}} V_h$, de telle façon que $V$ s'identifie à l'ens\supplied{emble} des $(x_i, x_j, x_h)$ tels que
\[
x_i \wedge x_j \wedge x_h = \varepsilon
\]
Si les $G_i$ sont \struck{commutatifs, ce qui est \uncertain{le cas} si $\ldots$ et \ill{}}, \add{commutatifs} donc (\uncertain{vu} l'iso\supplied{morphisme} \ill{}) isomorphes, les \uncertain{autres} $G$ sont et on a un \uncertain{système} transitif d'isomorphismes. On est ramené à la situation précédente, dès qu'on a donné une « origine » $o \in V$\ldots

\page{40}
\note{Numéro de l'auteur : 3. La page est d'une écriture rapide ; la prose de liaison est souvent incertaine.}
\emph{Remarque.} Le groupe des automorphismes de la situation donnée se divise en
\[
\mathfrak{S}_I \ (\simeq \mathfrak{S}_3) \qquad \text{\struck{$\mathrm{Aut}\, G_i \times \mathrm{Aut}\, G_i$}} \qquad \text{et}
\]
\[
\underbrace{\prod_{i \in I}{}_{\Gamma}\, \mathrm{Aut}(G_i)}_{\simeq\, \Gamma \text{ si les } G_i \text{ \uncertain{commutatifs}}}, \qquad \text{où} \quad
\Gamma \simeq \underbrace{\mathrm{Aut}\,\mathrm{Ext}(G_i)}_{\text{indépendant de } i \text{ à isom. près}}
\]
\note{Le $\prod$ porte en indice un $\Gamma$ appuyé (produit fibré au-dessus de $\Gamma$). Entre les deux, une tentative biffée, en partie noircie.}
\[
\prod G_i, \quad \text{opérant sur } (V_i) \text{ via} \quad g_j x_i g_h^{-1},\quad g_h x_j g_i^{-1},\quad g_i x_h g_j^{-1}
\]
\marginal{À droite, rattaché à $\Gamma$ par un trait : Il est \uncertain{plus} joli de distinguer $\mathfrak{S}_I$, $\mathrm{Aut}\,\mathrm{Ext}(G_i)$ et $G^I/Z$ (\uncertain{envoyé} diagonalement \uncertain{dans} $Z^I \subset G^I$)}
\marginal{En marge gauche, la feuille tournée, une note en grande partie surchargée : « \uncertain{produit} \uncertain{avec} \uncertain{conjugués}, \uncertain{il} faut \ill{} \struck{\ill{}} $Z^I$ / diagonale \ill{} $Z$ \ill{} \ill{} \ill{} »}

(c) Soit maintenant $I$ un \uncertain{ens\supplied{emble}} à quatre éléments, et \struck{\ill{}} donnons-nous encore d'abord les
\[
V \xrightarrow{\ \varphi_i\ } V_i \qquad i \in I
\]
dans $\mathcal{C}$, tels que $\forall J \subset I$, $\operatorname{card} J = 2$, on ait $V \xrightarrow{\ \varphi_J\ } V_J$ iso\supplied{morphisme}.

Pour toute \struck{subdivision} numérotation
\[
I = \{ i_1, i_2, i_3, i_4 \}
\]
on va définir une flèche
\[
u_{(i_1, i_2), (i_3, i_4)} : V_{i_1} \xrightarrow{\ \sim\ } V_{i_2} \tag{1}
\]
\note{À droite de (1), deux mots dont le premier est noirci (« \ill{} par \ill{} »).}
C'est l'iso\supplied{morphisme} canon\supplied{ique} défini en a), quand on fait abstraction de $\varphi_{i_4}$.
\note{La page s'arrête là ; l'argument se poursuit au-delà du lot.}

\end{document}
